\documentclass[11pt]{article}

\usepackage[margin=1in]{geometry}
\usepackage{amsmath, amssymb, mathtools}
\usepackage{enumitem}
\usepackage{parskip}
\usepackage{fancyhdr}
\usepackage[most]{tcolorbox}
\usepackage{xcolor}

% Page style
\pagestyle{fancy}
\fancyhf{}
\fancyhead[L]{CEC 315 -- Signals and Systems}
\fancyhead[R]{Spring 2026}
\fancyfoot[C]{\thepage}
\renewcommand{\headrulewidth}{0.4pt}

% Problem title box
\definecolor{examblue}{RGB}{30, 90, 170}
\newtcolorbox{problembox}[1]{
    colback=examblue!10,
    colframe=examblue!80!black,
    coltitle=white,
    fonttitle=\bfseries,
    title=#1,
    sharp corners,
    boxrule=0.6pt,
    left=6pt, right=6pt, top=4pt, bottom=4pt,
    breakable
}

\begin{document}

% -------- Title block --------
\begin{center}
    {\LARGE \textbf{CEC 315 -- Signals and Systems}}\\[2pt]
    Exam: Chapters 3 (Mock) \quad (Exam 3)
\end{center}

\noindent Rogelio Gracia Otalvaro \hfill Spring 2026

\vspace{4pt}
\hrule
\vspace{8pt}

\noindent\textbf{Name:} \rule{5cm}{0.4pt} \hfill \textbf{Date:} \rule{4cm}{0.4pt}

\vspace{6pt}

\noindent\textbf{Instructions:}
\begin{itemize}[leftmargin=*, itemsep=1pt, topsep=2pt]
    \item Total time: \textbf{60 minutes}. Total points: \textbf{100}.
    \item Show all work clearly. Box your final answers.
    \item Make sure your answers are clear and easy to understand.
    \item Read the text carefully and understand the questions, there is plenty of time.
    \item Assume all systems are causal and stable unless stated otherwise.
\end{itemize}

\noindent\textbf{Useful Formulas:}
\begin{align*}
&\text{Laplace pairs: } u(t)\leftrightarrow \tfrac{1}{s},\quad e^{-at}u(t)\leftrightarrow \tfrac{1}{s+a},\quad te^{-at}u(t)\leftrightarrow \tfrac{1}{(s+a)^2}\\
&\cos(\omega_0 t)u(t)\leftrightarrow \tfrac{s}{s^2+\omega_0^2},\quad \sin(\omega_0 t)u(t)\leftrightarrow \tfrac{\omega_0}{s^2+\omega_0^2}\\
&\text{Unilateral Laplace: } \mathcal{L}\{y'\} = sY - y(0^-),\quad \mathcal{L}\{y''\} = s^2Y - sy(0^-) - y'(0^-)\\
&\text{Z pairs: } a^n u[n]\leftrightarrow \tfrac{1}{1-az^{-1}},\ |z|>|a|;\quad -a^n u[-n-1]\leftrightarrow \tfrac{1}{1-az^{-1}},\ |z|<|a|\\
&(n+1)a^n u[n]\leftrightarrow \tfrac{1}{(1-az^{-1})^2}\\
&\text{Unilateral Z shift: } \mathcal{Z}\{y[n-1]\} = z^{-1}Y + y[-1]\\
&\mathcal{Z}\{y[n-2]\} = z^{-2}Y + z^{-1}y[-1] + y[-2]\\
&\text{Sampling theorem: } \omega_s > 2\omega_M\\
&\text{Closed loop (neg. feedback): } Q(s) = \dfrac{G(s)}{1+G(s)H(s)}
\end{align*}

\hrule
\vspace{10pt}

% -------- Part I: MC --------
\section*{Part I: Multiple Choice \hfill (5 questions $\times$ 4 points = 20 points)}

Circle the single best answer for each question.

\begin{enumerate}[leftmargin=*, label=\textbf{\arabic*.}]

\item A continuous-time signal has Fourier transform $X(j\omega)=0$ for $|\omega|>1500\pi$ rad/s. Which of the following sampling periods $T$ guarantees perfect reconstruction from the samples?
\begin{enumerate}[label=(\alph*)]
    \item $T = 1.0$ ms
    \item $T = 0.5$ ms
    \item $T = 2.0$ ms
    \item $T = 0.75$ ms
\end{enumerate}

\item A signal $x(t) = \cos(600\pi t)$ is sampled at $f_s = 400$ Hz. The reconstructed signal after an ideal lowpass filter with cutoff $f_s/2$ is a cosine of frequency:
\begin{enumerate}[label=(\alph*)]
    \item $300$ Hz
    \item $200$ Hz
    \item $100$ Hz
    \item $0$ Hz (DC)
\end{enumerate}

\item A causal LTI system has transfer function $H(s) = \dfrac{s+4}{(s+1)(s-2)}$. The system is:
\begin{enumerate}[label=(\alph*)]
    \item Stable, because it has a zero in the LHP.
    \item Stable, because two of its three singularities lie in the LHP.
    \item Unstable, because one pole is in the RHP.
    \item Marginally stable, because $|{-2}|<5$.
\end{enumerate}

\item The open-loop transfer function of a unity-feedback system is $G(s) = \dfrac{K}{s(s+4)}$. For what range of $K$ is the closed-loop system stable?
\begin{enumerate}[label=(\alph*)]
    \item $K < 0$
    \item $K > 0$
    \item $0 < K < 4$
    \item All real $K$ (it is always stable)
\end{enumerate}

\item A discrete-time signal has Z-transform $X(z)$ with poles at $z=0.3$ and $z=1.5$, and the signal is known to be two-sided (neither purely causal nor purely anti-causal). Which ROC is correct?
\begin{enumerate}[label=(\alph*)]
    \item $|z| > 1.5$
    \item $|z| < 0.3$
    \item $0.3 < |z| < 1.5$
    \item $|z| > 0.3$
\end{enumerate}

\end{enumerate}

\newpage

% -------- Part II --------
\section*{Part II: Problems \hfill (4 problems = 80 points)}

\begin{problembox}{Problem 1: Laplace Transform, ROC, and Inverse via PFE \hfill (20 points)}
\end{problembox}

\vspace{4pt}
Consider the Laplace transform
\[
X(s) \;=\; \frac{7s+2}{(s+2)(s-1)}, \qquad \text{ROC: } -2 < \operatorname{Re}\{s\} < 1.
\]
\begin{enumerate}[label=(\alph*)]
    \item \textbf{(4 pts)} Locate the poles of $X(s)$ in the $s$-plane and sketch the ROC.
    \item \textbf{(8 pts)} Perform a partial fraction expansion of $X(s)$. Show your cover-up work.
    \item \textbf{(6 pts)} Use the given ROC to determine whether each term corresponds to a right-sided or left-sided signal, and write the time-domain signal $x(t)$.
    \item \textbf{(2 pts)} Does the bilateral Fourier transform $X(j\omega)$ exist? Justify in one sentence.
\end{enumerate}

\newpage

\begin{problembox}{Problem 2: Unilateral Laplace -- Second-Order IVP with ZSR/ZIR \hfill (20 points)}
\end{problembox}

\vspace{4pt}
A causal LTI system is governed by
\[
\frac{d^2 y}{dt^2} \;+\; 4\,\frac{dy}{dt} \;+\; 3\,y(t) \;=\; 2\,u(t),
\]
with initial conditions $y(0^-) = 2$ and $y'(0^-) = -1$.
\begin{enumerate}[label=(\alph*)]
    \item \textbf{(4 pts)} Take the unilateral Laplace transform of both sides and solve algebraically for $Y(s)$ as a single rational function.
    \item \textbf{(4 pts)} Decompose $Y(s) = Y_{\text{zir}}(s) + Y_{\text{zsr}}(s)$, where $Y_{\text{zir}}(s)$ is the zero-input response (input set to zero, ICs kept) and $Y_{\text{zsr}}(s)$ is the zero-state response (ICs set to zero, input kept).
    \item \textbf{(8 pts)} Invert both pieces via partial fractions and write $y_{\text{zir}}(t)$ and $y_{\text{zsr}}(t)$ as explicit functions of $t$ for $t \ge 0$.
    \item \textbf{(4 pts)} Write the total response $y(t)$ and verify that it satisfies the initial conditions $y(0^-)=2$ and $y'(0^-)=-1$. Also state $y(\infty)$.
\end{enumerate}

\newpage

\begin{problembox}{Problem 3: Z-Transform, ROC, and Inverse with Mixed ROC \hfill (20 points)}
\end{problembox}

\vspace{4pt}
A discrete-time signal $x[n]$ has Z-transform
\[
X(z) \;=\; \frac{1 + z^{-1}}{\bigl(1 - 0.4\,z^{-1}\bigr)\bigl(1 + 0.2\,z^{-1}\bigr)}, \qquad \text{ROC: } 0.2 < |z| < 0.4.
\]
\begin{enumerate}[label=(\alph*)]
    \item \textbf{(3 pts)} Locate the poles and the zero. Sketch the pole-zero plot and shade the ROC on the complex $z$-plane.
    \item \textbf{(9 pts)} Perform a partial fraction expansion in $z^{-1}$. Show your cover-up calculations.
    \item \textbf{(6 pts)} Use the given annular ROC to decide whether each pole contributes a right-sided or left-sided term. Write $x[n]$ explicitly.
    \item \textbf{(2 pts)} Is $x[n]$ absolutely summable? Justify in one sentence using the ROC.
\end{enumerate}

\newpage

\begin{problembox}{Problem 4: Unilateral Z -- Difference Equation with Initial Conditions \hfill (20 points)}
\end{problembox}

\vspace{4pt}
A causal discrete-time system satisfies
\[
y[n] \;-\; \tfrac{1}{6}\,y[n-1] \;-\; \tfrac{1}{6}\,y[n-2] \;=\; x[n],
\]
with input $x[n] = u[n]$ and initial conditions $y[-1] = 6$, $y[-2] = 0$.
\begin{enumerate}[label=(\alph*)]
    \item \textbf{(4 pts)} Apply the unilateral Z-transform to the difference equation, carefully accounting for the initial conditions. Solve algebraically for $Y(z)$.
    \item \textbf{(4 pts)} Factor the characteristic polynomial and identify the system poles in the $z$-plane. Is the system stable?
    \item \textbf{(10 pts)} Perform a partial fraction expansion and invert to obtain $y[n]$ for $n \ge 0$ as an explicit closed-form expression.
    \item \textbf{(2 pts)} Verify your answer by computing $y[0]$ and $y[1]$ directly from the difference equation and confirming that they match your closed-form $y[n]$.
\end{enumerate}

\vfill
\begin{center}
    \textit{End of Exam --- Mock Exam 3 --- CEC 315, Spring 2026.}
\end{center}

\end{document}
